\chapter*{Conclusion et apprentissages personnels}
\label{chap:conclusion}
\markboth{\MakeUppercase{Conclusion et apprentissages personnels}}{}%
\addcontentsline{toc}{chapter}{Conclusion et apprentissages personnels}

Ce stage chez Dazzlifai m'a permis de travailler sur un sujet concret reliant automatisation documentaire, base de donnees, integration low-code et intelligence artificielle generative. Le projet a commence par une question simple: comment reduire le travail manuel autour des factures et devis? Il a progressivement evolue vers une solution complete, capable d'extraire, structurer, rechercher, interroger et presenter les informations.

Sur le plan technique, j'ai appris qu'un agent RAG utile ne se resume pas a appeler un modele de langage. La qualite de la reponse depend de tout ce qui vient avant: modele de donnees, extraction, validation, embeddings, recherche, filtres, gestion des erreurs et grounding. Une erreur dans une seule brique peut rendre la reponse finale moins fiable.

J'ai aussi compris l'importance de l'integration. Power Automate, Dataverse, Power Apps, Postgres, Gemini et Python n'ont pas les memes formats, les memes limites ni les memes manieres de signaler les erreurs. Faire fonctionner l'ensemble demande autant de rigueur dans les details que dans l'architecture generale.

Enfin, ce stage a change ma facon de voir le developpement logiciel applique a l'IA. L'IA apporte une grande valeur lorsqu'elle est encadree par des donnees fiables, des controles explicites et une experience utilisateur claire. Elle ne remplace pas la conception; elle la rend encore plus importante. Cette experience m'a donne envie de continuer a explorer les systemes intelligents, mais avec une attention particuliere a la fiabilite, a la tracabilite et a l'utilite reelle pour l'utilisateur.
